\documentclass[12pt,a4paper]{article}

% -------------------- Packages --------------------
\usepackage[a4paper,margin=1in]{geometry}
\usepackage{setspace}
\usepackage{titlesec}
\usepackage{enumitem}
\usepackage{array}
\usepackage{hyperref}
\usepackage{fancyhdr}

% -------------------- Page Formatting --------------------
\onehalfspacing
\setlength{\parindent}{0pt}
\setlength{\parskip}{6pt}

\titleformat{\section}
{\large\bfseries}
{\thesection.}
{0.5em}
{}

\titleformat{\subsection}
{\normalsize\bfseries}
{\thesubsection}
{0.5em}
{}

\setlist[itemize]{
    leftmargin=1.5em,
    itemsep=3pt
}

% -------------------- Header / Footer --------------------
\pagestyle{fancy}
\fancyhf{}
\lhead{\textbf{Journal 3}}
\rhead{\textbf{Smart City Civic Complaint System}}
\cfoot{\thepage}

% =========================================================
\begin{document}
% =========================================================

\begin{center}
    {\Large \textbf{JOURNAL 3}}\\[8pt]
    {\large \textbf{Frontend Implementation}}\\[6pt]
    {\normalsize \textbf{Smart City Civic Complaint and Management System}}
\end{center}

\vspace{12pt}

\begin{tabular}{|p{0.28\textwidth}|p{0.62\textwidth}|}
\hline
\textbf{Project} &
Smart City Civic Complaint and Management System \\ \hline

\textbf{Development Stage} &
Frontend Prototype Implementation \\ \hline

\textbf{Technology Used} &
React, Vite, React Router, JavaScript, CSS \\ \hline

\textbf{State Management} &
React Context API \\ \hline


\textbf{Student} &
Achintya Maurya \\ \hline
\end{tabular}

% =========================================================
\section{Objective}
% =========================================================

The objective of this development phase was to implement the frontend
prototype of the Smart City Civic Complaint and Management System using
React. The frontend was developed to provide an interactive interface for
citizens and department staff and to demonstrate the main complaint
management workflow before backend and database integration.

The work focused on application setup, navigation, role-based interfaces,
complaint submission and tracking, shared complaint state, officer
processing, and frontend integration.

% =========================================================
\section{Development Environment and Technologies}
% =========================================================

The frontend was developed using the following technologies:

\begin{itemize}
    \item \textbf{React} -- used for developing reusable user interface
    components.
    
    \item \textbf{Vite} -- used for creating and running the React
    development project.
    
    \item \textbf{JavaScript} -- used for frontend application logic.
    
    \item \textbf{React Router} -- used for navigation between different
    application pages.
    
    \item \textbf{CSS} -- used for styling dashboards, forms, cards,
    navigation bars, buttons, and status indicators.
    
    \item \textbf{React Context API} -- used for maintaining shared complaint
    information between frontend components.
\end{itemize}

% =========================================================
\section{Frontend Project Setup}
% =========================================================

The React frontend project was created using Vite. The project structure was
organized into separate folders for pages and application context so that the
frontend could be developed in a modular manner.

The main frontend structure includes separate page groups for:

\begin{itemize}
    \item Citizen
    \item Department Officer
    \item Field Worker
    \item Administrator
\end{itemize}

The development server was configured and the application was successfully
tested in the browser during implementation.

% =========================================================
\section{Application Routing}
% =========================================================

React Router was integrated into the application to provide navigation
between the different frontend interfaces.

The following major routes were implemented:

\begin{itemize}
    \item Login
    \item Register
    \item Citizen Dashboard
    \item Submit Complaint
    \item My Complaints
    \item Complaint Details
    \item Officer Dashboard
    \item Complaint Review
    \item Worker Dashboard
    \item Worker Task Details
    \item Administrator Dashboard
\end{itemize}

The routing structure allows each role to access its corresponding frontend
interface without creating separate applications.

% =========================================================
\section{Login and Registration Interface}
% =========================================================

A login interface was developed as the entry point of the application. The
interface contains fields for email and password along with a role-selection
field.

The available prototype roles are:

\begin{itemize}
    \item Citizen
    \item Department Officer
    \item Field Worker
\end{itemize}

The selected role determines which dashboard is opened after login.

A registration page was also implemented and connected to the login page.

At this stage, authentication is implemented only as a frontend prototype.
Actual authentication and authorization will be handled during backend
integration.

% =========================================================
\section{Citizen Dashboard}
% =========================================================

The Citizen Dashboard was developed as the main interface for citizens.
It provides access to the major complaint-related operations.

The dashboard contains:

\begin{itemize}
    \item Complaint statistics.
    \item Recent complaint information.
    \item Navigation to submit a complaint.
    \item Navigation to view submitted complaints.
    \item Complaint tracking access.
    \item Logout navigation.
\end{itemize}

The dashboard was designed to provide a simple interface through which a
citizen can submit and monitor civic complaints.

% =========================================================
\section{Complaint Submission}
% =========================================================

A dedicated \textit{Submit Complaint} page was implemented for reporting
civic issues.

The complaint form contains the following fields:

\begin{itemize}
    \item Complaint Category
    \item Complaint Description
    \item Location
    \item Photograph
\end{itemize}

The available categories in the frontend prototype include:

\begin{itemize}
    \item Garbage / Waste
    \item Street Light
    \item Road Damage
    \item Water Supply
    \item Drainage
    \item Other
\end{itemize}

The citizen enters the description of the civic issue and provides the
location where the issue has occurred. A photograph can also be selected
through the image upload field.

After submission, a temporary complaint ID is generated and a success
message is displayed to the citizen.

% =========================================================
\section{My Complaints}
% =========================================================

The \textit{My Complaints} page was implemented to allow citizens to view
their submitted complaints.

Each complaint card displays:

\begin{itemize}
    \item Complaint ID
    \item Complaint title
    \item Complaint category
    \item Location
    \item Submission date
    \item Current complaint status
\end{itemize}

A \textit{View Details} option is provided for opening the detailed
information of an individual complaint.

% =========================================================
\section{Complaint Details}
% =========================================================

The \textit{Complaint Details} page was implemented to display detailed
information about a complaint selected by the citizen.

The page contains:

\begin{itemize}
    \item Complaint title
    \item Complaint ID
    \item Category
    \item Location
    \item Submission date
    \item Current status
    \item Complaint description
\end{itemize}

A progress section was also included to represent the different stages of
complaint processing.

% =========================================================
\section{Shared Complaint State}
% =========================================================

To allow different frontend roles to work with common complaint information,
a shared state mechanism was implemented using the React Context API.

A file named \texttt{ComplaintContext.jsx} was created for maintaining the
complaint collection.

The complaint object stores information such as:

\begin{itemize}
    \item Complaint ID
    \item Title
    \item Category
    \item Location
    \item Date
    \item Citizen information
    \item Email
    \item Description
    \item Status
    \item Assigned worker
    \item Work progress
    \item Work notes
    \item Evidence
    \item Citizen verification status
\end{itemize}

Two main functions were implemented:

\begin{itemize}
    \item \texttt{addComplaint()} -- adds a new complaint to the shared
    complaint collection.
    
    \item \texttt{updateComplaint()} -- updates information of an existing
    complaint.
\end{itemize}

The \texttt{ComplaintProvider} was connected to the main React application
so that the complaint state could be accessed by different components.

% =========================================================
\section{Officer Dashboard Integration}
% =========================================================

The Officer Dashboard was connected with the shared complaint context.

Instead of using a separate static complaint list, the dashboard reads
complaint information from the common complaint state.

Dynamic statistics were implemented for:

\begin{itemize}
    \item Total Complaints
    \item Pending Verification
    \item In Progress
    \item Resolved
\end{itemize}

The dashboard also displays complaint information and provides a review
option for each complaint.

% =========================================================
\section{Complaint Review and Verification}
% =========================================================

The Complaint Review page was integrated with the shared complaint state.
The complaint ID from the selected route is used to identify the complaint
that needs to be reviewed.

The officer can view:

\begin{itemize}
    \item Complaint title
    \item Complaint ID
    \item Category
    \item Location
    \item Submission date
    \item Citizen name
    \item Citizen email
    \item Complaint description
    \item Current status
\end{itemize}

The following officer actions were implemented:

\begin{itemize}
    \item \textbf{Verify Complaint}
    \item \textbf{Reject Complaint}
    \item \textbf{Assign Field Worker}
\end{itemize}

When the officer verifies a complaint, its status is updated in the shared
frontend state. If the complaint is rejected, its status is changed to
\textit{Rejected}.

After verification, the officer can select a field worker and assign the
complaint.

% =========================================================
\section{Worker Assignment Integration}
% =========================================================

The frontend was prepared so that a verified complaint can be assigned to a
field worker.

The officer can select a worker from the available worker list. Once the
assignment is confirmed:

\begin{itemize}
    \item The assigned worker is stored with the complaint.
    \item The complaint status changes to \textit{Assigned}.
    \item The worker can access the corresponding task through the worker
    interface.
\end{itemize}

This integration connects the officer-side processing with the field-worker
workflow through the shared frontend complaint state.

% =========================================================
\section{Frontend Integration and Testing}
% =========================================================

The different frontend components were tested together to verify that the
application works as a single system.

The following activities were tested:

\begin{itemize}
    \item Starting the React application through the Vite development
    server.
    \item Navigation between different application routes.
    \item Role-based dashboard navigation.
    \item Citizen complaint form interaction.
    \item Complaint listing and detail navigation.
    \item Officer complaint review.
    \item Complaint verification and rejection.
    \item Field-worker assignment.
    \item Updating complaint information through React Context.
\end{itemize}

During integration, a duplicate router configuration caused a routing error.
This was resolved by maintaining a single \texttt{BrowserRouter} in
\texttt{main.jsx} and keeping the application routes inside \texttt{App.jsx}.

% =========================================================
\section{Current Frontend Result}
% =========================================================

The frontend prototype has reached a functional multi-role stage. The
Citizen, Department Officer, Field Worker, and Administrator interfaces have
been implemented as separate frontend views.

The complaint information can be shared between the implemented interfaces
through React Context, allowing the officer-side actions to update the
common complaint state.

The frontend currently demonstrates the application workflow without
requiring a backend server or database.

% =========================================================
\section{Current Limitations}
% =========================================================

The current frontend implementation has the following limitations:

\begin{itemize}
    \item Complaint information is currently maintained in React Context and
    is reset when the application is refreshed.
    
    \item Login is currently a frontend prototype and does not perform real
    authentication.
    
    \item Photograph and evidence inputs are currently handled at the
    frontend level.
    
    \item Persistent storage has not yet been implemented.
    
    \item Backend API integration is not part of the current frontend stage.
\end{itemize}

% =========================================================
\section{Learning Outcomes}
% =========================================================

The frontend implementation provided practical understanding of:

\begin{itemize}
    \item React component-based development.
    \item React Router and multi-page navigation.
    \item Role-based frontend organization.
    \item Interactive form handling in React.
    \item Shared state management using React Context API.
    \item Dynamic rendering based on complaint status.
    \item Frontend integration and debugging.
    \item Building an interactive prototype before backend integration.
\end{itemize}

% =========================================================
\section{Conclusion}
% =========================================================

The frontend implementation successfully established the user-facing
structure of the Smart City Civic Complaint and Management System. The major
citizen and officer interfaces, complaint handling functions, routing, and
shared complaint state were implemented and tested.

The completed frontend provides the foundation for further integration with
the remaining application modules and, subsequently, the backend API and
database.

% =========================================================
\end{document}
% =========================================================